\section{Background and Related Work}
\label{bg}
\subsection{Generative Recommendation}

In recent years, generative recommendation has attracted considerable attention in both academia and industry. This emerging paradigm, usually built on the Transformer architecture, reformulates recommendation as an end-to-end next-item generation task and thereby raises the performance ceiling of recommender systems.
Early work TIGER \citep{TIGER} employs residual quantization (RQ-VAE) \citep{rqvae} to convert text embeddings—extracted from an item’s title and description—into discrete semantic IDs, which are then used for next-item prediction. HSTU \citep{HSTU} propose a new architecture designed for high cardinality, non-stationary streaming recommendation data.
LC-Rec \citep{LCRec} aligns an LLM with semantic IDs via multi-task learning, enabling the model to “understand” the IDs and perform generative recommendation. Other studies investigate how to build better semantic IDs to enhance generation quality: RecForest \citep{recforest} applies hierarchical k-means clustering and treats the cluster indices as tokens, while EAGER \citep{eagar} and TokenRec \citep{TOKENRec} integrate semantic and collaborative signals directly into the tokenizer.

Very recently, generative recommendation has been rolled out at an industrial scale to address the drawbacks of traditional cascade systems. MTGR \citep{mtgr} keeps the original deep learning recommendation model (DLRM) features, introduces user-level compression, and speeds up both training and inference for large-scale deployment. OneRec \citep{OneRec} lowers the serving cost with a Lazy Decoder-Only Architecture and stabilises training through an improved reinforcement learning algorithm.

\subsection{LLM and RL}
Reinforcement learning (RL) trains an agent through repeated interaction with an environment so as to maximise cumulative return \citep{RLsurvey,RLsurvey2}. Within large-language-model (LLM) fine-tuning, RL with Human Feedback (RLHF) has become the de-facto recipe: it usually adopts Proximal Policy Optimisation (PPO) \citep{PPO} to align model behaviours with human preferences \citep{RLHF}. Unfortunately, PPO is memory-hungry at the billion-parameter scale, motivating a series of lighter alternatives. Direct Preference Optimisation (DPO) \citep{DPO} removes the value network and directly maximises the log-likelihood gap between preferred and dispreferred outputs; s-DPO \citep{sdpo} adapts this idea to recommendation by casting softmax negative mining as a pairwise-preference signal. Yet, preference-based methods remain off-policy and often plateau below on-line RL. Group-Relative Policy Optimisation (GRPO) \citep{deepseekmath} mitigates memory cost by normalising rewards inside a small group of roll-outs and replaces a learned reward model with rule-based heuristics, achieving strong gains on reasoning-heavy tasks such as mathematics and programming \citep{Deepseek-r1,o1}. 
% Very recent work continues to extend RL fine-tuning to broader domains—from mathematical proof search to general web search and information-retrieval scenarios \citep{R1-searcher,rm-r1,search-r1}.

Recent studies have begun to explore how SFT and RL jointly shape LLMs for generative recommendation.  
\citep{SFT-then-RL} argues that the two stages are complementary: a prolonged SFT phase first pushes accuracy to its limit, after which on-line RL with GRPO further compresses the token budget at inference time.  
\citep{RLvsSFT} shows that RL can largely recover the out-of-distribution accuracy lost during SFT by cancelling the directional drift of singular vectors rather than by discovering entirely new solutions.  
In contrast, \citep{RLReallyLearn} points out that the reasoning ability obtained through RL with verifiable rewards (RLVR) is bounded by the base model, whereas SFT can introduce genuinely new reasoning patterns, suggesting the need for more powerful RL paradigms such as continual scaling.  
\citep{RevisitingRL} adds a domain perspective: areas frequently encountered during pre-training (e.g., mathematics and code) profit from cross-domain RL, while low-exposure domains (e.g., logic and simulation) require in-domain RL for meaningful gains.  
\citep{EchoRL} observes that popular RL algorithms tend to converge to a single dominant output distribution, amplifying patterns already present in the pre-training data, yet they still display cross-task generalization.